% SPDX-License-Identifier: MIT-0
\documentclass[11pt,a4paper]{article}
\usepackage[T1]{fontenc}
\usepackage[utf8]{inputenc}
\usepackage{lmodern}
\usepackage[margin=25mm,headheight=14pt]{geometry}
\usepackage{amsmath,amssymb,amsthm,mathtools}
\usepackage{microtype,booktabs,array,longtable,enumitem}
\usepackage[dvipsnames]{xcolor}
\usepackage{fancyhdr}
\usepackage[colorlinks=true,linkcolor=MidnightBlue,citecolor=MidnightBlue,urlcolor=MidnightBlue]{hyperref}
\usepackage[nameinlink,noabbrev]{cleveref}
\hypersetup{pdftitle={Monotone Blocks and Reverse Replies: An Elementary Proof of the Large-Order Permutation-Avoidance Conjecture},pdfauthor={AI-assisted research report},pdfsubject={Permutation patterns; impartial games; a proposed solution of Ulfarsson's Conjecture 2.9}}
\pagestyle{fancy}\fancyhf{}
\fancyhead[L]{\small Monotone blocks and reverse replies}
\fancyhead[R]{\small Research draft}
\fancyfoot[C]{\thepage}
\setlength{\parskip}{0.3em}
\emergencystretch=2em
\setlist{itemsep=0.2em,topsep=0.4em}
\newtheorem{theorem}{Theorem}[section]
\newtheorem{lemma}[theorem]{Lemma}
\newtheorem{proposition}[theorem]{Proposition}
\newtheorem{corollary}[theorem]{Corollary}
\theoremstyle{definition}
\newtheorem{definition}[theorem]{Definition}
\newtheorem{example}[theorem]{Example}
\theoremstyle{remark}
\newtheorem{remark}[theorem]{Remark}
\DeclareMathOperator{\Av}{Av}
\DeclareMathOperator{\Sh}{Sh}
\DeclareMathOperator{\st}{st}
\DeclareMathOperator{\sg}{sg}
\DeclareMathOperator{\mex}{mex}
\newcommand{\N}{\mathbb N}
\newcommand{\Nk}{\mathcal N_k}
\newcommand{\id}{\iota}
\newcommand{\decr}{\delta}
\newcommand{\Tm}{\mathcal T}
\newcommand{\Um}{\mathcal U}
\newcommand{\replybound}{R(k)}
\newcommand{\shn}[1]{\Sh_k^{\mathrm{nm}}(#1)}
\newcommand{\code}[1]{\texttt{\detokenize{#1}}}

\title{\textbf{Monotone Blocks and Reverse Replies}\\[0.4em]
\large An elementary proof of the large-order\\permutation-avoidance conjecture}
\author{AI-assisted research report\\\small Prepared for Vladimir Reshetnikov}
\date{20 September 2026}

\begin{document}
\maketitle
\begin{abstract}
In the normal-play permutation-avoidance game, a move forbids a length-$k$
pattern that still occurs in a surviving permutation of length $n$. Ulfarsson's
Conjecture~2.9 asks whether, for each $k\ge3$ and all sufficiently large $n$,
the second player wins by replying to every pattern with its reversal.
We give a self-contained elementary proof with the explicit sufficient bound
\[
 n\ge \replybound:=k+(k+1)^2(k-2)^2+1.
\]
The main ingredient is a minimal-support lemma: in a family consisting of
singleton entries and one inflatable monotone interval, deleting singleton
entries can isolate exactly one member of any prescribed set of nonmonotone
patterns of a common length. A grid pigeonhole argument and the
Erd\H{o}s--Szekeres theorem extract such a family from every sufficiently long
legal witness. Reversal then supplies a witness for the required reply at the
\emph{same} length $n$. We also obtain an eventual pattern-isolation theorem,
a complement-reply corollary, and explicit small certificates. Exact Python
checks and two worked examples accompany the proof; no computation is a
premise of the universal result. The bound is not asserted to be optimal.
\end{abstract}

\noindent\fcolorbox{MidnightBlue}{blue!3}{\begin{minipage}{0.94\linewidth}
\small\textbf{Status of the result.} This is a proposed solution with a complete
proof, not a report of an independently accepted theorem. The argument has
been internally audited and tested on exact finite examples, but has not been
independently refereed or checked in a proof assistant. The source problem is
Conjecture~2.9 of arXiv:2603.16004v1, as retrieved on 20 September 2026.
\end{minipage}}

\setcounter{tocdepth}{1}
\begingroup\small
\setlength{\parskip}{0pt}
\makeatletter
\renewcommand{\l@section}[2]{\@dottedtocline{1}{0em}{1.8em}{#1}{#2}}
\makeatother
\tableofcontents
\endgroup
\newpage

\section{The problem and the main result}
\label{sec:problem}

\subsection{A move removes a pattern, not an entry}
Let $S_n$ be the permutations of $\{1,\ldots,n\}$. For a sequence of distinct
numbers $a_1,\ldots,a_j$, its \emph{standardization} $\st(a_1\cdots a_j)$ is
the unique permutation of length $j$ with the same pairwise comparisons.
A permutation $\pi$ contains $p\in S_k$ if
\[
 \st\bigl(\pi(i_1)\cdots\pi(i_k)\bigr)=p
 \quad\text{for some }i_1<\cdots<i_k.
\]
Write $p\le\pi$ for this relation. Otherwise $\pi$ avoids $p$. For a set of
patterns $F$, put
\[
 \Av_n(F)=\{\pi\in S_n:\ f\nleq\pi\text{ for every }f\in F\}.
\]

The game studied here is Ulfarsson's PAP game~\cite{Ulfarsson}. Its initial
position is $S_n$. When the already forbidden patterns are $F\subseteq S_k$,
a legal move chooses $p\in S_k$ occurring in some member of $\Av_n(F)$ and
changes the position to $\Av_n(F\cup\{p\})$. The player making the last legal
move wins. A selected pattern cannot be selected again, so play lasts at
most $k!$ moves.

For $p=p_1\cdots p_k$, let $p^r=p_k\cdots p_1$. The proposed strategy for
Player~II is simply
\[
 \boxed{\text{After Player I chooses }p,\text{ choose }p^r.}
\]
Although $p\ne p^r$ for $k\ge2$, that does not by itself make the reply legal:
all surviving occurrences of $p^r$ might have disappeared when $p$ was
forbidden. Proving that this cannot happen at sufficiently large $n$ is the
substantive issue.

\subsection{Published target and contribution}
Ulfarsson explicitly formulates eventual success of this strategy as
Conjecture~2.9~\cite{Ulfarsson}. His paper establishes the $k=3$ case and the
$k=4$ range $n\ge10$; the general statement is left open. The following
uniform theorem is the proposed resolution.

\begin{theorem}[Reverse-reply theorem]\label{thm:main}
For every integer $k\ge3$, define
\begin{equation}\label{eq:bound}
 \replybound=k+(k+1)^2(k-2)^2+1
 =k^4-2k^3-3k^2+5k+5.
\end{equation}
For every $n\ge\replybound$, Player~II wins PAP on $S_n$ with pattern length
$k$ by always replying with the reversal of Player~I's move. In particular,
$\sg(S_n,k)=0$ for every such $n$.
\end{theorem}

The proof does not classify the $k!$ patterns individually. It instead
minimizes the number of exceptional entries needed to retain \emph{either}
orientation of the proposed move. At a minimum, the two orientations cannot
both occur. This removes the need for a case-by-case list of companion
patterns.

\begin{center}
\begin{tabular}{@{}rrrrrrrrr@{}}
\toprule
$k$ &3&4&5&6&7&8&9&10\\
\midrule
$\replybound$ &20&105&330&791&1608&2925&4910&7755\\
\bottomrule
\end{tabular}
\end{center}

These are sufficient thresholds from one argument, not computed least
thresholds. In particular, the displayed $k=3,4$ bounds deliberately do not
compete with the sharper results already in the source. The new assertion is
uniformity in \emph{every} $k$.

\section{Elementary tools and the obstruction to naive pairing}
\label{sec:prelim}

Write
\[
 \id_k=12\cdots k,\qquad \decr_k=k\cdots21,\qquad
 \Nk=S_k\setminus\{\id_k,\decr_k\}.
\]
The elements of $\Nk$ are the nonmonotone length-$k$ patterns. Define the
$k$-shadow and its nonmonotone part by
\[
 \Sh_k(\pi)=\{p\in S_k:p\le\pi\},\qquad
 \shn{\pi}=\Sh_k(\pi)\cap\Nk.
\]
Classical containment is transitive, so if $\rho\le\pi$ then
$\Sh_k(\rho)\subseteq\Sh_k(\pi)$.

\begin{lemma}[Reversal compatibility]\label{lem:reversal}
For all $p\in S_k$ and $\pi\in S_n$,
\[
 p\le\pi\quad\Longleftrightarrow\quad p^r\le\pi^r.
\]
Consequently, if $F^r:=\{f^r:f\in F\}=F$, then $\Av_n(F)$ is closed under
reversal.
\end{lemma}
\begin{proof}
Reverse the order of the selected positions of an occurrence. Their values
are then read in reverse order, which standardizes to $p^r$. Reversal is an
involution, giving the converse. Apply the equivalence to each $f\in F$.
\end{proof}

\begin{remark}[Why symmetry alone is insufficient]
The lemma proves that $p$ and $p^r$ are either both legal or both illegal
\emph{before} $p$ is played from a reversal-closed state. A legal reply needs
more: a surviving permutation containing $p^r$ \emph{and avoiding $p$}.
Reversing a witness that contains both orientations still contains both and
therefore does not solve the problem.
\end{remark}

We use the following standard monotone-subsequence theorem, including its
elementary proof to keep all dependencies visible. Its classical source is
Erd\H{o}s and Szekeres~\cite{ES}.

\begin{lemma}[Erd\H{o}s--Szekeres]\label{lem:es}
If a permutation has more than $(a-1)(b-1)$ entries, it contains an
increasing subsequence of length $a$ or a decreasing subsequence of length
$b$. In particular, more than $(h-1)^2$ entries force a monotone subsequence
of length $h$.
\end{lemma}
\begin{proof}
At each position $j$, let $u_j$ and $v_j$ be the lengths of the longest
increasing and decreasing subsequences ending there. For $i<j$, if
$\pi(i)<\pi(j)$ then $u_j\ge u_i+1$; otherwise $v_j\ge v_i+1$. Thus all pairs
$(u_j,v_j)$ are distinct. If the two claimed subsequences are absent, all
pairs lie in $\{1,\ldots,a-1\}\times\{1,\ldots,b-1\}$, which has only
$(a-1)(b-1)$ elements.
\end{proof}

A \emph{monotone interval} of a permutation is a set of entries occupying
consecutive positions and consecutive values, with those entries increasing
or decreasing. In the extraction below, the interval property is required
only in a selected subpermutation. The corresponding points need not form
an interval in the original, much larger permutation.

\section{One monotone block and its stable shadow}
\label{sec:templates}

\begin{definition}[Pointed template]
A pointed template $\Tm=(\beta,j,\varepsilon)$ consists of a permutation
$\beta\in S_{m+1}$, an index $j$, and a sign $\varepsilon\in\{+,-\}$. For
$t\ge1$, the permutation $\Tm_t$ is obtained by replacing the entry
$\beta(j)$ by a monotone interval of length $t$ and sign $\varepsilon$.
Every other entry is retained as a singleton. Thus $|\Tm_t|=m+t$.
\end{definition}

Explicitly, set $v=\beta(j)$. Replace $v$ by
\[
 v,v+1,\ldots,v+t-1\quad\text{or}\quad v+t-1,\ldots,v+1,v,
\]
and replace every other value $x>v$ by $x+t-1$, leaving values $x<v$
unchanged. The sign specifies the choice of order. The block is consequently
in one position gap and one value gap relative to all singleton entries.

A subset $A$ of the singleton entries determines a subtemplate $\Tm^A$:
delete the other singleton entries from the skeleton, retain the marked
entry, and standardize. For every $t\ge1$,
\begin{equation}\label{eq:subtemplate}
 (\Tm^A)_t\le\Tm_t.
\end{equation}

\begin{lemma}[Stable shadows]\label{lem:stable}
For every pointed template and every $k\ge2$,
\begin{align}
 \Sh_k(\Tm_t)&=\Sh_k(\Tm_k) &&(t\ge k),\label{eq:stablefull}\\
 \shn{\Tm_t}&=\shn{\Tm_{k-1}} &&(t\ge k-1).\label{eq:stablenonmono}
\end{align}
\end{lemma}
\begin{proof}
A $k$-subsequence chooses some particular singleton entries and a number
$b$ of block entries. After the singleton choices and $b$ are fixed, its
pattern is fixed: all selected block entries have the same comparisons
with every selected singleton, and their internal order is specified by
$\varepsilon$. The choice of \emph{which} $b$ block entries is irrelevant.

For the full shadow, $0\le b\le k$, and all choices of $b$ that are possible
for $t\ge k$ are already possible for $t=k$. This proves
\eqref{eq:stablefull}. For a nonmonotone pattern, not all $k$ entries can
come from the monotone block. Thus at least one singleton is selected and
$b\le k-1$. All such choices are already available at $t=k-1$, proving
\eqref{eq:stablenonmono}.
\end{proof}

Denote this stable nonmonotone shadow by
\[
 H_k(\Tm):=\shn{\Tm_{k-1}}.
\]
It can be computed from a permutation of length $m+k-1$, no matter how large
the eventual inflated permutation will be. By \eqref{eq:subtemplate},
\begin{equation}\label{eq:shadowmonotonicity}
 A\subseteq A'\quad\Longrightarrow\quad
 H_k(\Tm^A)\subseteq H_k(\Tm^{A'}).
\end{equation}

\begin{remark}[The missing monotone pattern is harmless, but not absent]
The full shadow need not stabilize at $k-1$: a block with $k-1$ entries
cannot by itself supply its length-$k$ monotone pattern, whereas a block with
$k$ entries can. We will inflate only while all forbidden patterns are
nonmonotone. This is exactly why the smaller stabilization length $k-1$ is
valid in the final bound.
\end{remark}

\section{The minimal-core isolation lemma}
\label{sec:minimal}

The next lemma is the decisive step. It concerns arbitrary sets of targets,
not just two reversals.

\begin{lemma}[Minimal-core isolation]\label{lem:minimal}
Let $Q\subseteq\Nk$ be nonempty, and suppose $H_k(\Tm)\cap Q\ne\varnothing$.
There is a subtemplate $\Um$ obtained by deleting singleton entries such that
\begin{equation}\label{eq:isolate}
 |H_k(\Um)\cap Q|=1.
\end{equation}
Moreover, $\Um$ has at most $k$ singleton entries, and at least one.
\end{lemma}
\begin{proof}
Choose an inclusion-minimal subset $A$ of the singleton entries such that
\[
 H_k(\Tm^A)\cap Q\ne\varnothing,
\]
and put $\Um=\Tm^A$. Such an $A$ exists because the original finite set of
singletons has the required property. Let $m=|A|$.

Consider any occurrence of any member of $Q$ in $\Um_{k-1}$. That occurrence
must use \emph{every} singleton of $A$. Indeed, if it omitted a singleton
$a\in A$, the same occurrence would survive in
$(\Tm^{A\setminus\{a\}})_{k-1}$, contradicting the minimality of $A$.

It follows that every occurrence of any target uses exactly the same $m$
singleton entries, and exactly $k-m$ entries of the monotone block. But these
choices determine a \emph{unique} relative order: the singleton entries are
fixed, all block entries lie in the same position and value gaps relative
to them, and the internal order of the block is fixed. Thus all target
occurrences standardize to one and the same permutation. Since at least one
target occurs, \eqref{eq:isolate} follows.

The existing occurrence uses all $m$ singletons, so $m\le k$. Also $m\ne0$:
a monotone block alone cannot contain a nonmonotone target. This proves the
last assertion.
\end{proof}

\begin{remark}[What is, and is not, being minimized]
We minimize the singleton support needed to retain \emph{some} member of
$Q$. Minimizing the support of one preselected target would not justify
excluding the other targets: another target could use a proper subset of
that support. Allowing the retained target to change during deletion is
essential. In the game application, reversal symmetry subsequently orients
the isolated target in the desired direction.
\end{remark}

\begin{example}[The smallest illustrative pair]\label{ex:small}
Let $k=3$, $Q=\{132,231\}$, and inflate the first entry of $1342$ increasingly.
The resulting family is
\[
 \Tm_t=1,2,\ldots,t,\ t+2,\ t+3,\ t+1.
\]
For $t\ge2$, its nonmonotone $3$-shadow is $\{132,231\}$. Delete the
singleton $t+2$ and standardize. The smaller family is
\[
 \Um_t=1,2,\ldots,t,\ t+2,\ t+1,
\]
whose only nonmonotone $3$-pattern is $132$. Its reversal has only $231$.
For instance, at block length two the successive permutations are
\[
 12453\quad\longrightarrow\quad1243
 \quad\longrightarrow\quad3421.
\]
The middle and last permutations respectively isolate the two orientations.
\end{example}

\subsection{A finite, terminating minimization algorithm}
The proof gives a direct procedure. Starting with $\Tm$, delete any singleton
whose deletion leaves at least one target in the stable shadow. Repeat until
no deletion is possible. Every successful step reduces the number of
singletons by one, so the process terminates. At termination, the same
omitted-singleton argument proves that exactly one target remains.

There is no uniqueness assertion for the resulting template or isolated
target. Different deletion orders can give different certificates. This
freedom is useful: a reply proof needs one suitable witness, not a canonical
one.

\section{Extracting an inflatable witness from a long permutation}
\label{sec:extract}

\begin{lemma}[Same-cell extraction]\label{lem:extract}
Suppose $\pi$ contains a marked occurrence of $q\in S_k$, and let $h\ge2$.
If
\begin{equation}\label{eq:extractionbound}
 |\pi|\ge k+(k+1)^2(h-1)^2+1,
\end{equation}
then $\pi$ has a subpermutation $\Tm_h$ formed from that occurrence as $k$
singleton entries and a monotone block of length $h$ in a pointed template.
In particular, $q\le\Tm_h\le\pi$.
\end{lemma}
\begin{proof}
Plot the entries as points $(i,\pi(i))$. The $k$ marked positions partition
the unmarked positions into $k+1$ gaps, including the two exterior gaps.
Likewise, the $k$ marked values partition the unmarked values into $k+1$
gaps. The unmarked points therefore lie in $(k+1)^2$ open cells, indexed by
their position gap and value gap.

There are at least $(k+1)^2(h-1)^2+1$ unmarked points. Some cell contains
at least $(h-1)^2+1$ points. Read its values from left to right and apply
\cref{lem:es}; this yields an increasing or decreasing subsequence of length
$h$ entirely inside the cell.

Retain exactly those $h$ points and the $k$ marked points, and standardize.
The $h$ retained cell points are consecutive in position among the selected
points, because no marked point lies inside their position gap. They are
also consecutive in value, because no marked value lies inside their value
gap. Thus they form a monotone interval relative to the marked singletons.
Collapsing this interval to one marked entry gives the required template.
\end{proof}

\begin{corollary}[Stable extraction]\label{cor:extract}
Let $k\ge3$ and $|\pi|\ge\replybound$. If $q\in\Nk$ occurs in $\pi$, there is
a template $\Tm$ with exactly $k$ singletons such that
\[
 q\in H_k(\Tm)\subseteq\shn{\pi}.
\]
Consequently, if $\pi$ avoids $F\subseteq\Nk$, then $\Tm_t$ avoids $F$ for
every $t\ge k-1$.
\end{corollary}
\begin{proof}
Use $h=k-1$ in \cref{lem:extract}. The extracted permutation is
$\Tm_{k-1}\le\pi$, and the original occurrence of $q$ uses only its
singletons. The shadow inclusion follows by containment transitivity.
Inflating beyond $k-1$ preserves the nonmonotone shadow by
\cref{lem:stable}; all members of $F$ lie in that part of the shadow.
\end{proof}

We can now combine extraction and minimization without mentioning a game.

\begin{theorem}[Eventual pattern isolation]\label{thm:isolation}
Let $k\ge3$, let $F,Q\subseteq\Nk$, and suppose some
$\pi\in\Av_n(F)$ with $n\ge\replybound$ contains a member of $Q$.
There exist a template $\Um$ with $1\le m\le k$ singleton entries and a
pattern $q_*\in Q$ such that, for every $t\ge k-1$,
\begin{equation}\label{eq:universal-isolation}
 \Um_t\in\Av_{m+t}(F),\qquad
 \Sh_k(\Um_t)\cap Q=\{q_*\}.
\end{equation}
In particular, for every $N\ge2k-1$ there is a permutation of length $N$
avoiding $F$ and containing exactly one member of $Q$.
\end{theorem}
\begin{proof}
Choose a member of $Q$ occurring in $\pi$ and apply \cref{cor:extract}.
Its template $\Tm$ has $H_k(\Tm)\cap Q\ne\varnothing$ and
$H_k(\Tm)\cap F=\varnothing$. Apply \cref{lem:minimal} to obtain $\Um$.
Deleting singletons can only shrink the stable shadow, so $\Um$ still avoids
$F$. The target intersection is a singleton, say $\{q_*\}$, and stability
makes these statements valid for every $t\ge k-1$.

For the last assertion choose $t=N-m$. Since $m\le k$ and $N\ge2k-1$,
we have $t\ge k-1$. The resulting permutation has exactly $N$ entries.
\end{proof}

\begin{remark}[The quantifiers matter]
The hypothesis requires one witness at a length at least $\replybound$.
The conclusion then supplies isolated-target witnesses at every length
$N\ge2k-1$. It does \emph{not} say that any witness of length $2k-1$ is
already sufficient for the extraction, nor does it assert that a prescribed
member of $Q$ can always be isolated. The latter orientation is supplied by
symmetry only in the next section.
\end{remark}

\section{Legal reverse replies and the winning strategy}
\label{sec:gameproof}

\begin{theorem}[Same-length reverse witness]\label{thm:reply}
Let $k\ge3$, $n\ge\replybound$, and $F\subseteq\Nk$ satisfy $F^r=F$.
If $p\in\Nk$ is legal from $\Av_n(F)$, then there exists
\[
 \rho\in\Av_n(F\cup\{p\})\quad\text{such that}\quad p^r\le\rho.
\]
Thus $p^r$ remains legal immediately after $p$ is played.
\end{theorem}
\begin{proof}
A legal move has a witness $\pi\in\Av_n(F)$ containing $p$.
Apply \cref{thm:isolation} with $Q=\{p,p^r\}$, and choose the resulting
isolating permutation $\sigma$ to have length \emph{exactly} $n$.
This choice is permitted because $n\ge\replybound\ge2k-1$.

The two targets are distinct: equality $p=p^r$ would force the first and
last entries of a permutation of length at least two to agree. If $\sigma$
contains $p^r$, it avoids $p$ and is the desired $\rho$. Otherwise it contains
$p$ and avoids $p^r$. Set $\rho=\sigma^r$. By \cref{lem:reversal}, $\rho$
contains $p^r$, avoids $p$, and still avoids $F$ because $F^r=F$.
\end{proof}

For completeness, the monotone endgame and the passage from replies to a
winning strategy are proved explicitly rather than imported.

\begin{lemma}[The monotone endgame]\label{lem:endgame}
If $F\subseteq\Nk$ and $n\ge k$, both $\id_k$ and $\decr_k$ are legal from
$\Av_n(F)$. After one is played, the other remains legal. If moreover
$n\ge(k-1)^2+1$, forbidding both leaves no surviving permutation.
\end{lemma}
\begin{proof}
The increasing permutation $\id_n$ and decreasing permutation $\decr_n$
avoid every nonmonotone pattern. They therefore survive $F$ and witness the
two monotone moves. If $\id_k$ is chosen, $\decr_n$ still survives and
contains $\decr_k$; the opposite case is identical with the orders reversed.
The final assertion is \cref{lem:es} with $a=b=k$.
\end{proof}

\begin{proof}[Proof of \cref{thm:main}]
At the beginning of a round, just before Player~I moves, maintain the
following invariant: either the game has ended on Player~II's previous
move, or the forbidden set $F$ is reversal-closed and contains no monotone
pattern. The invariant holds initially because $F=\varnothing$.

If Player~I chooses a nonmonotone $p$, \cref{thm:reply} makes $p^r$ a legal
reply. Adding both $p$ and $p^r$ preserves reversal-closure and introduces
neither monotone pattern, so the invariant is restored.

If Player~I chooses a monotone pattern, \cref{lem:endgame} makes its reversal
a legal reply. After that reply no permutation survives, since
$n\ge\replybound\ge(k-1)^2+1$. Thus Player~II wins immediately.

There is no possibility of an infinite run consisting only of nonmonotone
rounds. Each round forbids two distinct new patterns, and there are only
$k!-2$ nonmonotone patterns. Before a monotone round, both monotone moves
remain legal. Therefore play reaches a monotone round after finitely many
rounds, and Player~II makes the last legal move.

In a finite impartial normal-play game, a position is losing for the next
player exactly when its Sprague--Grundy value is zero; a proof of this
standard equivalence is included in \cref{app:sg}. Hence
$\sg(S_n,k)=0$ as asserted.
\end{proof}

\subsection{Stronger local conclusion}
The proof applies not only to the initial state $S_n$, but to every state
$\Av_n(F)$ with $F\subseteq\Nk$ and $F^r=F$. From any such state, the next
player loses against reversal replies by the other player when
$n\ge\replybound$. No reachability assumption on $F$ is needed.

It is important that the forbidden set at a round boundary contains no
monotone pattern. Once Player~I first chooses one, we do not invoke the
inflation argument again: the very next reverse reply finishes the game.

\section{A fully explicit length-330 certificate}
\label{sec:example330}

The following example exercises the first length-$5$ bound in
\cref{thm:main}. It is illustrative, not an additional assumption.
Let $p=12354$, so $p^r=45321$. Take the skeleton
\[
 \beta=124635
\]
and inflate its fifth entry, of value $3$, decreasingly. There are five
singleton entries, so block length $t=325$ gives
\begin{equation}\label{eq:bigpi}
 \pi=1,2,328,330,327,326,\ldots,3,329\in S_{330}.
\end{equation}
The singleton entries themselves, at positions $1,2,3,4,330$, form $p$.
The stable nonmonotone $5$-shadow of this template is
\begin{align*}
 H=\{&12354,12435,12453,12534,12543,13524,14325,\notag\\
     &14532,15324,15432,35214,43215,45321,53214\}.
\end{align*}
It contains both $p$ and $p^r$. Define the reversal-closed forbidden set
\[
 F=\mathcal N_5\setminus(H\cup H^r).
\]
Here $|F|=92$. All members of $F$ are avoided by \eqref{eq:bigpi}.

Minimization with $Q=\{12354,45321\}$ deletes three singleton entries.
The skeletons, with the marked entry underlined and its block always
\emph{decreasing}, are
\begin{equation}\label{eq:trace}
 1246\underline{3}5
 \quad\longrightarrow\quad
 135\underline{2}4
 \quad\longrightarrow\quad
 24\underline{1}3
 \quad\longrightarrow\quad
 23\underline{1}.
\end{equation}
At every step at least one target remains. The final template has two
singletons and, for every $s\ge4$, the permutation
\[
 \Um_s=s+1,s+2,s,s-1,\ldots,1.
\]
Its only nonmonotone $5$-pattern is $45321$: selecting both singletons and
three block entries yields that pattern, while using at most one singleton
yields a decreasing pattern. Thus the length-$330$ reply witness is
\begin{equation}\label{eq:rho330}
 \rho=329,330,328,327,\ldots,1.
\end{equation}
It contains $45321$, avoids $12354$, and avoids all $92$ patterns of $F$.
No final reversal happens to be needed in this example.

The file \code{data/worked_certificate.json} records the full input and
output permutations, the exact forbidden set, the selected cell, the
monotone subblock, every deletion, and the final stable shadow. All of its
pattern assertions can be checked on representatives with a block of length
four or six; enumerating the $\binom{330}{5}$ subsequences of the large
permutations is unnecessary.

\section{Constructive algorithm and certificate size}
\label{sec:algorithm}

\subsection{Input and output}
For a nonmonotone move, the constructive input consists of $k$, $n$, a
reversal-closed $F\subseteq\Nk$, and a witness $\pi\in\Av_n(F)$ together with
one marked occurrence of $p$. Supplying the occurrence avoids a separate
pattern-search problem. The output is an explicit $\rho\in S_n$ witnessing
the legality of $p^r$ after $p$.

\begin{center}
\begin{minipage}{0.94\linewidth}
\hrule\medskip
\textbf{Reverse-witness construction}
\begin{enumerate}[label=\arabic*.,leftmargin=1.6em]
\item Split the unmarked entries of $\pi$ into the $(k+1)^2$ cells defined
by the marked occurrence. Choose a cell with at least $(k-2)^2+1$ entries.
\item Find an increasing or decreasing subsequence of length $k-1$ in that
cell. Retain it together with the marked occurrence, producing $\Tm_{k-1}$.
\item Repeatedly delete a singleton whenever the stable shadow still meets
$Q=\{p,p^r\}$. Let $\Um$ be the resulting template with $m$ singletons.
\item Form $\sigma=\Um_{n-m}$. If its isolated target is $p^r$, output
$\sigma$; otherwise output $\sigma^r$.
\end{enumerate}
\medskip\hrule
\end{minipage}
\end{center}

\subsection{Finite proof certificates}
A certificate can store the original marked occurrence, the selected
$k-1$ cell points, the template and deletion trace, the isolated orientation,
and the final block length. The initial finite representative has $2k-1$
entries, and every later representative has at most that many.

The verifier checks that the original selected points standardize to the
stated template; checks finite pattern shadows of the representatives;
checks that the final shadow contains exactly one target and avoids $F$;
and, when necessary, invokes reversal-closure of $F$. The mathematical
stability lemma then transports these finite checks to the length-$n$
output. Thus the certificate is independent of $n$ except for the indices
of selected points and the final block length.

More precisely, apart from the supplied $F$ and permutation, it contains
$O(k^2)$ small skeleton entries if the whole deletion trace is retained,
and $O(k\log n)$ bits for the selected original indices. Retaining only the
initial and final templates reduces the structural record further, though
the full trace is easier to inspect.

\subsection{A conservative running-time bound}
With a marked occurrence supplied, gap assignment takes $O(nk)$ elementary
comparisons. A straightforward dynamic program finds the required monotone
subsequence in $O(n^2)$ time, or faster methods can be substituted without
changing the proof. Each finite shadow test enumerates at most
$\binom{2k-1}{k}$ subsequences; a direct standardization takes $O(k^2)$
comparisons. At most $k(k+1)$ deletion attempts suffice for a naive repeated
scan. For fixed $k$, the overall construction is therefore polynomial in
$n$; a conservative bound is
\[
 O\!\left(n^2+k^4\binom{2k-1}{k}\right)
\]
comparison operations, apart from representing and testing membership in
$F$. There is no assertion of polynomial dependence on $k$.

The supplied implementation instead computes stable shadows symbolically:
it chooses the singleton subset and only the \emph{number} of selected
block entries. An independent direct subsequence routine checks these
symbolic computations in the test suite.

\section{Extensions and the limits of the argument}
\label{sec:extensions}

\subsection{Complement replies}
For $p\in S_k$, define $p^c$ by $(p^c)_i=k+1-p_i$. Complementation preserves
classical containment, has no fixed permutation when $k\ge2$, and exchanges
the two monotone patterns.

\begin{corollary}[Complement strategy]\label{cor:complement}
For every $k\ge3$ and $n\ge\replybound$, Player~II also wins PAP on $S_n$ by
replying to $p$ with $p^c$.
\end{corollary}
\begin{proof}
At a round boundary keep $F$ complement-closed and free of monotone
patterns. Apply \cref{thm:isolation} to $Q=\{p,p^c\}$. If necessary,
complement the isolated witness to orient it toward $p^c$. It still avoids
$F$. The two monotone patterns are paired by complementation, so
\cref{lem:endgame} and the same termination argument apply.
\end{proof}

This is a direct consequence of the isolation mechanism, not a separate
claim that a new game has been solved. The mechanism requires a symmetry
with no fixed target and a suitable endgame. For example, reversal followed
by complementation can fix patterns, so this argument does not automatically
give a reply strategy for that operation.

\subsection{Forbidden patterns of bounded, unequal lengths}
The targets must have a common length in \cref{lem:minimal}, but the
forbidden patterns need not. The distinction permits the following useful
extension.

\begin{corollary}[Bounded forbidden lengths]\label{cor:mixedF}
Let $k\ge3$, let $Q\subseteq\Nk$, and let $F$ consist of nonmonotone
patterns of lengths at most $\ell$. Set $K=\max(k,\ell)$. If some
$F$-avoiding permutation $\pi$ containing a member of $Q$ has length at least
\[
 k+(k+1)^2(K-2)^2+1,
\]
then for every $N\ge k+K-1$ there is an $F$-avoiding permutation of length
$N$ containing exactly one member of $Q$.
\end{corollary}
\begin{proof}
Extract a same-cell monotone block of length $K-1$ around a marked target
occurrence. The nonmonotone shadow of every length $j\le K$ is already
stable at that block length. Hence enlarging the block preserves avoidance
of every member of $F$. Minimize singleton support while retaining a target
of length $k$. The proof of \cref{lem:minimal} is unchanged, and leaves at
most $k$ singletons. Inflating to total length $N\ge k+K-1$ leaves at least
$K-1$ block entries, preserving all required avoidance conditions.
\end{proof}

This is a statement about permutation classes and witnesses. It does not
settle a game allowing moves of different lengths: the pairing and endgame
would require a separate analysis.

\subsection{Why the main hypotheses are substantive}
The proof has four boundaries. First, the stable representative must be
long enough for every relevant nonmonotone pattern; replacing $k-1$ by a
smaller number without a separate argument could introduce new forbidden
patterns during inflation. Second, every member of $Q$ has the same length;
otherwise using all the singletons need not force the same number of block
entries. Third, avoiding monotone forbidden patterns is not preserved by
unbounded monotone inflation. Fourth, reversal-closure of $F$ is what
permits us to orient the isolated target without losing the previous
avoidance conditions.

None of these conditions is silently assumed away. They hold exactly where
they are used in the PAP proof.

\section{Exact computations and reproducibility}
\label{sec:computations}

The universal proof consists of \cref{lem:stable,lem:minimal,lem:extract}
and the game argument. The following computations were carried out as
independent finite consistency checks and as an aid to discovery.

\subsection{Two shadow implementations}
The function \code{shadow_direct} enumerates actual index subsets of a
finite permutation and standardizes their values. The function
\code{Template.stable_shadow} instead enumerates singleton subsets and the
number of block entries. These use different enumeration schemes, although
both necessarily use elementary rank standardization.

For every skeleton in $S_{k+1}$, every marked slot, and both block signs,
the two routines were compared at block lengths $k-1,k,k+1$ for $k=3,4$.
The resulting exhaustive counts are as follows.

\begin{center}
\begin{tabular}{@{}rrrr@{}}
\toprule
$k$ & Templates & Reverse-pair minimizations & Other target-set tests\\
\midrule
3&192&404&2192\\
4&1200&6482&0\\
\bottomrule
\end{tabular}
\end{center}

For each reverse-pair minimization, the final target intersection and its
minimality under every singleton deletion were checked by the direct
routine. At $k=3$, all nonempty target subsets were also tested whenever
present in the starting template.

There were additionally $120$ deterministically sampled templates for each
of $k=5,6,7$, with comparisons at block lengths $k-1$ and $k+1$, and
minimizations for target sets containing up to four sampled patterns and
their reversals. These are sampled tests, not exhaustive tests in those
three lengths.

\subsection{Constructed replies and extraction tests}
For each $k=3,4,5,6,7$, twenty full length-$\replybound$ replies were
constructed from one-block witnesses, using the maximal reversal-closed
forbidden set disjoint from each input shadow. This gives $100$ certified
large-permutation outputs. Their final compressed representatives were
checked directly for avoidance and the required target orientation.

A separate set of tests applied same-cell extraction to shuffled, unstructured
permutations with empty forbidden set. The suite also checked the sharp
$(h-1)^2$ monotone-subsequence obstruction for $2\le h\le9$, and twenty
shuffled sequences of length $(h-1)^2+1$ for each of those values.

Every assertion in the executed suite passed. The test seed is
$20260920$ and is explicitly separate from the operating-system randomness
used to choose the research area. The exact output is in
\code{data/verification.json}; the complete length-$330$ example is in
\code{data/worked_certificate.json}.

\subsection{The discovery computation}
Before the minimal-core lemma was recognized, a finite exploration enumerated
one-block shadows with exactly $k$ singletons. A reversal support records
which nonmonotone reversal pairs meet a shadow, without recording their
orientations. For each support $S$ arising in the enumeration and every
pair in $S$, the program looked for a separating template whose support
was contained in $S$.

\begin{center}
\begin{tabular}{@{}rrrrr@{}}
\toprule
$k$ & Templates & Distinct full shadows & Supports & Failed supports\\
\midrule
5&8640&7336&3502&0\\
6&70560&61816&30794&0\\
\bottomrule
\end{tabular}
\end{center}

Those tests suggested a support-minimization principle, but by themselves
would not prove the all-$k$ theorem. The later argument permits deleting
singletons and handles every $k$ without these tables. The package retains
a reproducible exploration script and summary output so that the discovery
record is distinguished from the proof dependencies.

\subsection{Running the package}
The main verification requires Python 3.10 or later and no third-party
packages. From the archive root, run
\begin{verbatim}
python code/verify.py
python code/explore_supports.py --k 5
python code/explore_supports.py --k 6
latexmk -pdf -interaction=nonstopmode -halt-on-error article.tex
\end{verbatim}
The exploration is optional. The first command checks the examples and
implementation; the last builds this article from its source. No network
connection is required to run the checks or build the PDF once the ordinary
\LaTeX{} packages are installed.

\section{What remains to be determined}
\label{sec:remaining}

The existence assertion in the selected conjecture is the part addressed
by \cref{thm:main}. The argument leaves the quantitative threshold largely
unoptimized. The quartic expression has a transparent origin: $k$ marked
entries create $(k+1)^2$ cells, and a cell with more than $(k-2)^2$ points
contains the required block. This is an upper bound on the worst-case
extraction cost, not evidence that quartic growth is necessary for the game.

One concrete next question is the least eventual reverse-strategy threshold
for $k=5$. Our result bounds it above by $330$, but neither the proof nor the
recorded computation determines it. More generally, improving the use of the
marked occurrence or exploiting restrictions imposed by $F$ might lower
the uniform bound.

The report does not compute the full Sprague--Grundy function, analyze
optimal play lengths, prove a mis\`ere analogue, or handle arbitrary
nonclassical pattern notions. The construction is tailored to classical
subsequence containment and normal play. These limitations do not leave a
missing case of Conjecture~2.9: they concern different or stronger questions.

\appendix
\section{A compact proof audit}
\label{app:audit}

The following checklist isolates the places where an incorrect proof of a
pairing strategy could otherwise hide a quantifier or length error.

\begin{longtable}{@{}p{0.30\linewidth}p{0.66\linewidth}@{}}
\toprule
\textbf{Potential gap} & \textbf{Where it is closed}\\
\midrule
\endhead
A reverse witness still contains the just-forbidden pattern.
& Minimal-core isolation gives \emph{exactly one} of $p,p^r$ before any
symmetry is applied.\\[0.4em]
The long monotone block is not an interval in the original permutation.
& Only the selected $2k-1$ points are retained; being in one cell makes the
block an interval in that selected subpermutation.\\[0.4em]
Inflating a witness creates a forbidden pattern.
& Every forbidden pattern at this stage is nonmonotone of length $k$, so its
occurrence is already visible at block length $k-1$.\\[0.4em]
The minimized occurrence uses a different singleton set.
& Minimality forces \emph{every} target occurrence to use all retained
singletons, not merely one chosen occurrence.\\[0.4em]
Different block choices give different target patterns.
& The block is monotone and in one position/value gap; fixing the number
chosen and all singletons fixes every relative comparison.\\[0.4em]
The isolated target is the wrong orientation.
& Reverse the witness; $F^r=F$ preserves prior avoidance.\\[0.4em]
The witness has the wrong length.
& With $m\le k$ singleton entries, choose block length $n-m\ge k-1$ to obtain
exactly $n$ entries.\\[0.4em]
A monotone move violates the inflation hypotheses.
& Such a move is handled only by the separate monotone endgame, which ends
the game on the very next reply.\\[0.4em]
The reply rule does not imply eventual termination.
& Every legal move forbids a new one of the $k!$ patterns. Before a monotone
round, both monotone moves remain legal.\\[0.4em]
A finite test is treated as a universal proof.
& No assertion in the main proof depends on the computational tables or on
an unverified classification.\\
\bottomrule
\end{longtable}

\section{Selection, exclusion, and source-status record}
\label{app:selection}

The research area was selected by one call to Python's
\code{secrets.randbelow(96)} after fixing a list of $96$ mathematical areas.
The recorded zero-based result was $9$, hence entry $10$:
\emph{Permutation patterns}. The draw was recorded at
\code{2026-09-20T21:58:38.955311+00:00}. The complete list and output are in
\code{area_selection.json}, and \code{select_area.py} refuses to overwrite an
existing draw. The record is an audit trail of the executed choice, not a
cryptographic proof of when a human first saw the list.

The entire attached \code{manifest(1).tex} was read, including its final
continuation. It contains no entry for PAP, reverse-reply strategies, or
Ulfarsson's Conjecture~2.9. The manifest's permutation-counting and
permutation-avoidance entries concern different problems; none is selected
here. A hash of the supplied manifest and the exclusion rationale are
recorded in \code{notes/selection_and_exclusions.md}.

One initial candidate from the stack-sorting literature was rejected after
finding a later paper proving it~\cite{LiEtAl}. The selected source was then
checked directly in its versioned arXiv HTML and PDF. Conjecture~2.9 appears
on PDF page~5, and the general-$k$ discussion appears in Section~6.1.
Targeted searches by paper identifier, title, and conjecture wording did not
identify a later resolution. This is a documented source check, not a
complete bibliographic guarantee that no independent or unpublished proof
exists.

The contribution claimed here is the elementary proof of the stated general
reverse-reply theorem through same-cell extraction and minimal-core
isolation. Classical monotone-subsequence arguments, the PAP rules, and the
previous small-$k$ results are not claimed as new. The precise priority of
the abstract isolation lemma beyond this application has not been settled
by an exhaustive literature search.

\section{Why the Sprague--Grundy conclusion follows}
\label{app:sg}

For a finite impartial normal-play game, define the Sprague--Grundy value
recursively by
\[
 \sg(X)=\mex\{\sg(Y):Y\text{ is a legal follower of }X\},
\]
where the minimal excludant is the least nonnegative integer not in the set.
A terminal position has value zero. Induct on the maximum remaining play
length. If $\sg(X)=0$, every follower has nonzero value and, by induction,
is winning for the next player; therefore $X$ is losing for its next player.
If $\sg(X)>0$, zero must occur among the follower values, so the next player
can move to a losing position and win. This proves the equivalence used in
\cref{thm:main}. PAP is finite because each selected pattern becomes
permanently illegal, so the induction applies.

\begin{thebibliography}{9}
\bibitem{Ulfarsson}
Henning Ulfarsson,
\emph{A Permutation Avoidance Game with Reverse Replies and Monotone Traps},
arXiv:2603.16004v1, 2026.
Conjecture~2.9; Section~6.1.
\url{https://arxiv.org/abs/2603.16004v1}.
Versioned HTML and PDF consulted 20 September 2026.

\bibitem{ES}
P. Erd\H{o}s and G. Szekeres,
\emph{A combinatorial problem in geometry},
Compositio Mathematica \textbf{2} (1935), 463--470.
\url{https://www.numdam.org/item/CM_1935__2__463_0/}.

\bibitem{LiEtAl}
Yang Li, Sergey Kitaev, Zhicong Lin, and Jing Liu,
\emph{A bijection between 321- and 213-avoiding permutations preserving
$t$-stack-sortability}, arXiv:2507.09187, 2025.
\url{https://arxiv.org/abs/2507.09187}.
Used only to exclude an already-resolved candidate, not as a proof dependency.
\end{thebibliography}

\end{document}
